\documentclass[10pt,a4paper]{amsart}
\usepackage[margin=19mm]{geometry}
\usepackage{amsmath,amssymb,mathtools}
\usepackage[scr=rsfs,cal=euler]{mathalfa}
\usepackage{xfrac}
\usepackage[unicode,hidelinks,bookmarks=false]{hyperref}
\newcommand{\ie}{i.e.\ }
\newcommand{\cf}{cf.\ }
\newcommand{\resp}{respectively\ }
\newcommand{\Subsection}{\subsection}
\providecommand{\nobreakdash}{\nobreak\hbox{-}}
\setlength{\emergencystretch}{2em}
\multlinegap0pt
\usepackage[shortlabels]{enumitem}
\usepackage{amssymb}
\usepackage{mathtools}
\usepackage{xcolor}
\usepackage{algorithm}\usepackage{algpseudocode}
\usepackage{tikz}
\newtheorem{assumption}{Assumption}
\newtheorem{theorem}{Theorem}
\title{From Microscopic Damage to Macroscopic Games: A Dimensionality Reduction of Stem Cell Homeostasis}
\author{Jiguang Yu, Louis Shuo Wang, Shihan Ban}
\date{Source: arXiv:2603.12713v1 (CC BY 4.0)}
\begin{document}
\maketitle

\noindent\emph{Source and licence.} From Microscopic Damage to Macroscopic Games: A Dimensionality Reduction of Stem Cell Homeostasis. Version arXiv:2603.12713v1 (CC BY 4.0), distributed by arXiv under the Creative Commons Attribution 4.0 International licence, http://creativecommons.org/licenses/by/4.0/, as recorded in the arXiv licence metadata for this exact version and shown on the abstract page https://arxiv.org/abs/2603.12713v1. Full licence notice: \texttt{licenses/arxiv-2603.12713-CC-BY-4.0.txt}.\par\smallskip

\noindent\textit{Editorial scope.} Counted unit: the article's theorem thm:balance-laws (source0.tex lines 287--307). The article's complete original proof of it is delivered with it (source0.tex lines 309--344, one \texttt{proof} environment of 36 lines). No mathematics is written, repaired or re-derived: the statement and the proof are the article's own lines, in the article's own notation, and no retained line is substituted or re-flowed.  Retained uncounted as the source-local context the proof needs: lines 179--186; lines 188--190; lines 192--194; lines 202--209; lines 239--249.  Nothing was omitted from any retained span, so the package carries no omission record.  No retained line contains a citation, so the wrapper carries no bibliography and no bibliography entry.  No retained line contains a figure environment, an image-inclusion command or drawing code, so no image is delivered and the package declares no asset dependency.  Editorial formatting only, in addition: an AMS \texttt{amsart} preamble that restates the article's own declarations and macros used by the retained lines, namely the environments \texttt{assumption}, \texttt{theorem} on separate counters, together with the standard packages those lines need.  The retained environments print in this document as Assumption 1, Theorem 1 in order of appearance, whereas the article prints the same items with the numbering of its own full document.  The complete native source of the article as distributed in the arXiv e-print for this exact version is bundled unchanged as \texttt{sources/196291/source0.tex}.
\begin{equation}\label{eq:PDE_model}
\begin{cases}
\partial_t P + v_P \partial_x P
= \mathcal B_P[P](t,x) - \lambda_P(\bar W(t))\,P(t,x) + \lambda_R(\bar P(t))\,W(t,x),\\[4pt]
\partial_t W + v_W \partial_x W
= \mathcal B_W[P](t,x) - \big(\delta(x)+\lambda_R(\bar P(t))\big)\,W(t,x),
\end{cases} \qquad x>0,\ t>0.
\end{equation}

\begin{equation}\label{eq:inflow_bc}
P(t,0)=W(t,0)=0,\qquad t\ge 0,
\end{equation}

\begin{equation}\label{eq:IC}
P(0,x)=P_0(x),\qquad W(0,x)=W_0(x),\qquad x>0.
\end{equation}

\begin{align}\label{eq:Hill_maps}
\begin{aligned}
p_1(\bar W)=\frac{\hat p_1}{1+(k_1\bar W)^{m_1}},\quad
p_2(\bar W)=\frac{\hat p_2}{1+(k_2\bar W)^{m_2}},\\[4pt]
\lambda_P(\bar W)=\frac{\hat\lambda_P}{1+(k_3\bar W)^{m_3}},\quad
\lambda_R(\bar P)=\frac{\hat\lambda_R}{1+(k_4\bar P)^{m_4}},
\end{aligned}
\end{align}

\begin{assumption}\label{ass:basic}
\begin{enumerate}[label=(H\arabic*), leftmargin=2.6em]
\item \label{H1}
The maps $p_1,p_2,\lambda_P,\lambda_R$ are bounded and globally Lipschitz on $[0,\infty)$
(e.g.\ the Hill families in \eqref{eq:Hill_maps} satisfy this).
\item \label{H2}
$\delta\in L^\infty([0,\infty))$ is measurable and nonnegative.
\item \label{H3}
$P_0,W_0\in L^1_+([0,\infty))$.
\end{enumerate}
\end{assumption}

\begin{theorem}[Balance laws for totals]
\label{thm:balance-laws}
Assume Assumption~\ref{ass:basic}. Then the PDE system \eqref{eq:PDE_model}--\eqref{eq:IC}
admits a unique global nonnegative solution
$(P,W)\in C([0,\infty);L^1([0,\infty))\times L^1([0,\infty)))$.
Moreover, the total masses satisfy, for a.e.\ $t>0$,
\begin{equation}\label{eq:balance_laws_general}
\begin{cases}
\displaystyle
\frac{d\bar P}{dt}
=\Big(p_1(\bar W)-p_2(\bar W)\Big)\lambda_P(\bar W)\,\bar P
+\lambda_R(\bar P)\,\bar W,\\[8pt]
\displaystyle
\frac{d\bar W}{dt}
=\Big(1-p_1(\bar W)+p_2(\bar W)\Big)\lambda_P(\bar W)\,\bar P
-\lambda_R(\bar P)\,\bar W
-\int_0^\infty \delta(x)\,W(t,x)\,dx,
\end{cases}
\end{equation}
with initial conditions $\bar P(0)=\|P_0\|_{1}$ and $\bar W(0)=\|W_0\|_{1}$.
\end{theorem}

\begin{proof}[Proof sketch (derivation of \eqref{eq:balance_laws_general})]
We outline the formal computation; the rigorous argument follows by testing the weak formulation
against cutoffs $\psi_R(x)$ and passing $R\to\infty$, where $\psi_R \equiv 1$ on $[0,R]$, $\operatorname{supp} \psi_R \subset [0,R+1]$ and $\lVert\psi_R' \rVert_{\infty}\le C$ with $C$ independent of $R$ (cf.\ standard transport weak formulations).

Integrate the $P$-equation in \eqref{eq:PDE_model} over $x\in(0,\infty)$:
\[
\frac{d}{dt}\int_0^\infty P\,dx
+\int_0^\infty v_P\partial_x P\,dx
=\int_0^\infty \mathcal B_P[P]\,dx
-\lambda_P(\bar W)\int_0^\infty P\,dx
+\lambda_R(\bar P)\int_0^\infty W\,dx.
\]
The transport term gives the boundary flux
$\int_0^\infty v_P\partial_x P\,dx = v_P P(t,\infty)-v_P P(t,0)$.
Under integrability and the inflow boundary condition \eqref{eq:inflow_bc}, this term vanishes:
$P(t,\infty)=0$ (in the sense of $L^1$ tails) and $P(t,0)=0$ (in the sense of traces for $L^1$ solutions to transport equations).
An identical computation holds for the $W$-equation.

It remains to compute $\int \mathcal B_P[P]\,dx$ and $\int \mathcal B_W[P]\,dx$.
Using the change of variables $y=x/\alpha$,
\[
\int_0^\infty \frac{1}{\alpha}P\!\left(t,\frac{x}{\alpha}\right)\,dx
=\int_0^\infty P(t,y)\,dy=\bar P(t),
\]
and similarly for $\alpha\in\{\alpha_i,\beta_i,\gamma_i\}$, together with $p_1+p_2+p_3=1$,
one obtains
\[
\int_0^\infty \mathcal B_P[P]\,dx=(2p_1(\bar W)+p_3(\bar W))\lambda_P(\bar W)\bar P
=\big(1+p_1(\bar W)-p_2(\bar W)\big)\lambda_P(\bar W)\bar P,
\]
\[
\int_0^\infty \mathcal B_W[P]\,dx=(2p_2(\bar W)+p_3(\bar W))\lambda_P(\bar W)\bar P
=\big(1-p_1(\bar W)+p_2(\bar W)\big)\lambda_P(\bar W)\bar P.
\]
Substituting these identities yields \eqref{eq:balance_laws_general}.
\end{proof}

\end{document}
